\textbf{结论名称：} 奇函数在原点有定义则必过原点（$f(0)=0$）。\\
\textbf{适用条件：} 函数 $f(x)$ 为奇函数，且定义域包含 $0$。\\
\textbf{变量说明：} $f:D\to\mathbb{R}$，满足 $f(-x)=-f(x)$，且 $0\in D$。\\
\textbf{核心公式：} \[ \boxed{f(0)=0} \]\\
\textbf{使用场景：} 已知奇函数求参数；利用奇偶性简化求值；判断函数图像是否过原点。\\
\textbf{简单例子：} 已知定义在 $\mathbb{R}$ 上的函数 $f(x)=\dfrac{1}{2^x+1}+m$ 为奇函数，求 $m$。\\
由 $f(0)=0$ 得 $\dfrac{1}{2^0+1}+m=\dfrac{1}{2}+m=0$，解得 $m=-\dfrac12$。\\
（验证：此时 $f(x)=\dfrac{1}{2^x+1}-\dfrac12=\dfrac{1-2^x}{2(2^x+1)}$ 满足 $f(-x)=-f(x)$。）